(1) 利用 \([a^\dagger a,a]=-a\)、\([\sigma_z,\sigma_+]=2\sigma_+\)，有
\[
[N,a\sigma_+]=(-a)\sigma_++a\sigma_+=0,
\]
共轭项也为零，故 \([H,N]=0\)。在 \(\{|e,n\rangle,|g,n+1\rangle\}\) 中自由部分同为 \(\hbar\omega(n+1/2)\)，矩阵
\[
H_n=\hbar\omega(n+1/2)I+\hbar g\sqrt{n+1}\begin{pmatrix}0&1\\1&0\end{pmatrix}.
\]
(2) 演化为
\[
|\psi(t)\rangle=e^{-i\omega(n+1/2)t}\big[\cos(g\sqrt{n+1}t)|e,n\rangle-i\sin(g\sqrt{n+1}t)|g,n+1\rangle\big],
\]
故 \(P_e=\cos^2(g\sqrt{n+1}t)\)。对于 \(n=0\) 为真空Rabi振荡 \(\cos^2(gt)\)，两条dressed能级劈裂 \(2\hbar g\sqrt{n+1}\)。
(3) 相干态的Poisson权重 \(p_n=e^{-\bar n}\bar n^n/n!\)，
\[
P_e(t)=\tfrac12+\tfrac12 e^{-\bar n}\sum_{n=0}^\infty{\bar n^n\over n!}\cos(2gt\sqrt{n+1}).
\]
对窄相对光子数分布 \(n=\bar n+\delta n,\mathrm{Var}(n)=\bar n\)，相位线性化为 \(2gt\sqrt{\bar n+1}+(gt/\sqrt{\bar n+1})\delta n\)。其Poisson特征函数包络近似 \(\exp[-g^2t^2\bar n/(2(\bar n+1))]\simeq e^{-g^2t^2/2}\)，所以按 \(1/e\) 包络定义 \(t_{\mathrm{collapse}}\simeq\sqrt2/g\)。邻近n相位差 \(2gt(\sqrt{n+2}-\sqrt{n+1})\simeq gt/\sqrt{\bar n}\)，约等于 \(2\pi\) 时初次复活，\(t_{\mathrm{rev}}\simeq2\pi\sqrt{\bar n}/g\)。高阶色散使复活并非精确完美。
(4) 初态 \(\langle H\rangle=\hbar\omega(\bar n+1/2)\)，相互作用项初期为零；因H不显含时间，\(\langle H\rangle\) 全程恒定。失谐将Rabi频率改为 \(\sqrt{\Delta^2+4g^2(n+1)}\)，损耗需要主方程；当 \(g/\omega\) 不再很小时反旋转项影响激发数守恒，需使用量子Rabi模型。
